\documentclass[11pt]{article}
\usepackage{mathblatt}


\begin{document}
\pfheftstil
\pfheftkopf{Daten}{P10}
\pfrueckblick
\begin{pfaufg}{}{1.}{}
Ordne der Größe nach: 52,2; 19,2; 85,1; 35,3.
\fusshilfe{\mbox{1:~19,2 < 35,3 < 52,2 < 85,1}}
\end{pfaufg}
\begin{pfaufg}{}{2.}{}
Berechne: 1400 : 7; 528 : 11; 16 : 4.
\rechenraster{1}
\fusshilfe{\mbox{2:~200; 48; 4}}
\end{pfaufg}
\begin{pfaufg}{}{3.}{}
Ergänze die Tabelle Anteil | Winkel.
\par\smallskip\noindent \begingroup\renewcommand{\arraystretch}{1.3}\begin{tabular}{|c|c|}\hline \textbf{Anteil} & \textbf{Winkel} \\ \hline 360° & \rule{0pt}{3ex}\hspace{1.6cm} \\ \hline 180° & \rule{0pt}{3ex}\hspace{1.6cm} \\ \hline 90° & \rule{0pt}{3ex}\hspace{1.6cm} \\ \hline 36° & \rule{0pt}{3ex}\hspace{1.6cm} \\ \hline 3,6° & \rule{0pt}{3ex}\hspace{1.6cm} \\ \hline\end{tabular}\endgroup\par
\fusshilfe{\mbox{3:~360°; 180°; 90°; 36°; 3,6°}}
\end{pfaufg}
\begin{pfaufg}{}{4.}{}
Eine Achse reicht mit 8 Kästchen bis 160. Wie viel ist ein Kästchen? Wie viele Kästchen hoch ist eine Säule für 90?
\rechenraster{1}
\fusshilfe{\mbox{4:~20 je Kästchen; 4,5 Kästchen}}
\end{pfaufg}
\begin{pfaufg}{}{5.}{}
Ergänze die Tabelle Wort | Prozent für: die Hälfte, ein Viertel, drei Viertel, jede 15.
\par\smallskip\noindent \begingroup\renewcommand{\arraystretch}{1.3}\begin{tabular}{|c|c|}\hline \textbf{Wort} & \textbf{Prozent} \\ \hline 50 \% & \rule{0pt}{3ex}\hspace{1.6cm} \\ \hline 25 \% & \rule{0pt}{3ex}\hspace{1.6cm} \\ \hline 75 \% & \rule{0pt}{3ex}\hspace{1.6cm} \\ \hline $\tfrac{1}{15}$ & \rule{0pt}{3ex}\hspace{1.6cm} \\ \hline\end{tabular}\endgroup\par
\fusshilfe{\mbox{5:~50 \%; 25 \%; 75 \%; $\tfrac{1}{15}$ = $\tfrac{20}{3}$ \% $\approx$ 6,7 \%}}
\end{pfaufg}
\pfstufekopf[sminimummaximumspannw]{Minimum, Maximum, Spannweite}{in 5 der letzten 5 Prüfungen}
\pfgruppe{}
\begin{pfaufg}{}{6.}{}
Jemand fragt: „Welche Schuhgröße kommt am häufigsten vor?“ Kreuze an, welche Kenngröße gefragt ist.
\pfkreuzzeile{Spannweite}\pfkreuzzeile{Modalwert}\pfkreuzzeile{Median}\pfkreuzzeile{Mittelwert}
\fusshilfe{\mbox{6:~Modalwert}}
\end{pfaufg}
\pfgruppe{}
\begin{pfaufg}{}{7.}{}
Die Punkte sind: $14$, $9$, $22$, $17$, $11$. Gib das Minimum (kleinster Wert) und das Maximum (größter Wert) an.
\par\noindent Min = \leerfeld , Max = \leerfeld \par
\fusshilfe{\mbox{7:~$22$}}
\end{pfaufg}
\begin{pfaufg}{NI ’22}{8.}{}
In Hannover fallen im Januar im Mittel 52 mm Regen, im Juni 74 mm. Berechne den Unterschied.
\rechenraster{1}
\fusshilfe{\mbox{8:~74 mm $-$ 52 mm = 22 mm}}
\end{pfaufg}
\begin{pfaufg}{NRW ’25}{9.}{}
Eine App zeigt Jakubs Bildschirmzeit von Montag bis Donnerstag. Gib die größte Bildschirmzeit an.
\par\smallskip\noindent \fbox{\parbox{0.9\linewidth}{\footnotesize Skizze: Säulendiagramm Bildschirmzeit: Mo 1 h 10 min; Di 2 h 6 min; Mi 2 h 43 min; Do 49 min; Fr leer}}\par
\fusshilfe{\mbox{9:~2 h 43 min (Mittwoch)}}
\end{pfaufg}
\pfgruppe{}
\begin{pfaufg}{BY ’24}{10.}{}
Die Spieler einer Handballmannschaft haben die Schuhgrößen 41, 43, 45, 43, 44 und 45. Ergänze die Schuhgröße von Spieler 7 so, dass die Spannweite 5 beträgt.
\par\smallskip\noindent \sachtabelle{l}{Spieler-Nr. 1 bis 7}{Schuhgröße 41, 43, 45, 43, 44, 45, leer}\par
\fusshilfe{\mbox{10:~40 oder 46}}
\end{pfaufg}
\begin{pfaufg}{NI ’22}{11.}{}
Die Tabelle zeigt die mittleren Monatstemperaturen auf dem Brocken im Jahr 2020. Nenne den wärmsten und den kältesten Monat.
\par\smallskip\noindent \sachtabelle{l}{Monat Jan bis Dez}{Temperatur in °C: $-$4, $-$3, $-$2, 2, 4, 9, 12, 15, 10, 4, 3, 1}\par
\fusshilfe{\mbox{11:~kältester Monat: Januar ($-$4 °C)}}
\end{pfaufg}
\begin{pfaufg}{P10 ’14}{12.}{2 BE}
Eine Tankstelle hat im Sommer 2012 zehn Wochen lang einmal pro Woche die Preise für die Kraftstoffe E 10 und Diesel notiert (in Cent pro Liter). Lies aus der Tabelle den niedrigsten und den höchsten Preis für E 10 ab (Minimum und Maximum).
\par\smallskip\noindent \sachtabelle{lcc}{Datum & E 10 & Diesel}{03.07. & 154,2 & 139,6\\ 10.07. & 157,4 & 142,5\\ 17.07. & 158,8 & 144,2\\ 24.07. & 159,7 & 145,4\\ 31.07. & 159,1 & 146,2\\ 07.08. & 161,0 & 147,1\\ 14.08. & 164,8 & 150,7\\ 21.08. & 169,2 & 154,0\\ 28.08. & 168,1 & 153,1\\ 04.09. & 167,0 & 152,0}\par
\fusshilfe{\mbox{12:~Minimum 154,2 ct; Maximum 169,2 ct}}
\end{pfaufg}
\begin{pfaufg}{P10 ’14}{13.}{2 BE}
Eine Tankstelle hat im Sommer 2012 zehn Wochen lang einmal pro Woche die Preise für die Kraftstoffe E 10 und Diesel notiert (in Cent pro Liter). Berechne die Spannweite der Dieselpreise in diesen zehn Wochen.
\par\smallskip\noindent \sachtabelle{lcc}{Datum & E 10 & Diesel}{03.07. & 154,2 & 139,6\\ 10.07. & 157,4 & 142,5\\ 17.07. & 158,8 & 144,2\\ 24.07. & 159,7 & 145,4\\ 31.07. & 159,1 & 146,2\\ 07.08. & 161,0 & 147,1\\ 14.08. & 164,8 & 150,7\\ 21.08. & 169,2 & 154,0\\ 28.08. & 168,1 & 153,1\\ 04.09. & 167,0 & 152,0}\par
\fusshilfe{\mbox{13:~14,4 ct}}
\end{pfaufg}
\begin{pfaufg}{P10 ’22}{14.}{3 BE}
Von 2017 bis 2021 wurden jedes Jahr 1 200 Jugendliche gefragt, wie oft sie Bücher lesen. Das Säulendiagramm zeigt, wie viele von ihnen mehrmals pro Woche ein Buch lesen. Gib den kleinsten und den größten Wert im Diagramm an (Minimum und Maximum). Bestimme den Durchschnitt (das arithmetische Mittel) der fünf Werte.
\par\smallskip\noindent \begin{tikzpicture}\begin{axis}[ybar,width=8.5cm,height=4.6cm,ymin=390,ymax=508.75,symbolic x coords={2017,2018,2019,2020,2021},xtick=data,enlarge x limits=0.08,bar width=6mm,ylabel={Anzahl},ylabel style={font=\small},tick label style={font=\small},nodes near coords,every node near coord/.append style={font=\scriptsize,/pgf/number format/.cd,use comma,1000 sep={}},ymajorgrids,grid style={mbgitter}]\addplot[fill=mbkasten,draw=black] coordinates {(2017,468.0) (2018,432.0) (2019,485.0) (2020,470.0) (2021,400.0)};\end{axis}\end{tikzpicture}\par
\pfzwfrage{Berechne die Summe der fünf Werte.}
\rechenraster{2}
\fusshilfe{\mbox{14:~Minimum 400 (2021), Maximum 485 (2019); Mittel 451}}
\end{pfaufg}
\pfgruppe{}
\begin{pfaufg}{P10 ’19}{15.}{1 BE}
Gib die Spannweite dieser Daten an: 9,7 m; 9,9 m; 9,6 m; 9,5 m; 9,8 m; 9,7 m
\rechenraster{1}
\fusshilfe{\mbox{15:~0,4 m}}
\end{pfaufg}
\pfgruppe{}
\begin{pfaufg}{P10 ’24}{16.}{2 BE}
Das Säulendiagramm zeigt, wie viel Niederschlag 2021 in Brandenburg in jedem Monat gefallen ist (in mm). Ermittle die Spannweite der monatlichen Niederschläge im Jahr 2021 in mm.
\par\smallskip\noindent \begin{tikzpicture}\begin{axis}[ymin=0,ymax=100,ytick distance=20,tick label style={font=\small},xtick pos=bottom,ytick pos=left,yticklabel style={/pgf/number format/use comma,/pgf/number format/1000 sep={\,},/pgf/number format/fixed,/pgf/number format/precision=1},ymajorgrids,grid style={mbgitter},xtick={0,1,2,3,4,5,6,7,8,9,10,11},xticklabels={{Jan},{Feb},{Mär},{Apr},{Mai},{Jun},{Jul},{Aug},{Sep},{Okt},{Nov},{Dez}},xmin=-0.5,xmax=11.5,bar width=5mm,nodes near coords,every node near coord/.append style={font=\scriptsize,/pgf/number format/use comma,/pgf/number format/1000 sep={\,},/pgf/number format/fixed,/pgf/number format/precision=1},ybar,width=13.5cm,height=5.2cm,ylabel={Niederschlag in mm},ylabel style={font=\small}]\addplot[fill=mbkasten,draw=black] coordinates {(0,52.2) (1,39.6) (2,35.3) (3,19.2) (4,49.3) (5,85.1) (6,36.1) (7,82.4) (8,21.7) (9,54.6) (10,64.1) (11,42.1)};\end{axis}\end{tikzpicture}\par
\rechenraster{1}
\fusshilfe{\mbox{16:~65,9 mm}}
\end{pfaufg}
\begin{pfaufg}{P10 ’23}{17.}{3 BE}
Auf der Welt werden wohl über 7 000 Sprachen gesprochen. Die Tabelle nennt die fünf Sprachen mit den meisten Muttersprachlern (Anzahl in Millionen): Arabisch 290, Chinesisch 1 300, Englisch 500, Hindi 525, Spanisch 389. Gib für diese fünf Werte das Minimum, das Maximum und die Spannweite an.
\par\smallskip\noindent \sachtabelle{lr}{Sprache & Muttersprachler in Millionen}{Arabisch & 290\\ Chinesisch & 1 300\\ Englisch & 500\\ Hindi & 525\\ Spanisch & 389}\par
\fusshilfe{\mbox{17:~Minimum 290 Mio.; Maximum 1300 Mio.; Spannweite 1010 Mio.}}
\end{pfaufg}
\pfgruppe{}
\begin{pfaufg}{P10 ’20}{18.}{2 BE}
Das Säulendiagramm zeigt die Besucherzahlen eines Freibads an den Tagen der ersten Ferienwoche. Die beiden folgenden Aussagen sind richtig: „Die Besucherzahlen der ersten Ferienwoche haben eine Spannweite von 150.“ „Am Sonntag gab es ein Drittel mehr Besucher als am Mittwoch.“ Weise beide Aussagen rechnerisch nach.
\par\smallskip\noindent \begin{tikzpicture}\begin{axis}[ybar,width=11.1cm,height=4.6cm,ymin=0,ymax=336.00,symbolic x coords={Mo,Di,Mi,Do,Fr,Sa,So},xtick=data,enlarge x limits=0.08,bar width=6mm,ylabel={Besucher je Wochentag},ylabel style={font=\small},tick label style={font=\small},nodes near coords,every node near coord/.append style={font=\scriptsize,/pgf/number format/.cd,use comma,1000 sep={}},ymajorgrids,grid style={mbgitter}]\addplot[fill=mbkasten,draw=black] coordinates {(Mo,170.0) (Di,130.0) (Mi,210.0) (Do,180.0) (Fr,190.0) (Sa,240.0) (So,280.0)};\end{axis}\end{tikzpicture}\par
\pfzwfrage{Berechne ein Drittel der Besucherzahl vom Mittwoch.}
\rechenraster{3}
\fusshilfe{\mbox{18:~280 $-$ 130 = 150; 210 + 210 : 3 = 280}}
\end{pfaufg}
\pfgruppe{}
\begin{pfaufg}{P10 ’21}{19.}{5 BE}
Das Säulendiagramm gibt an, wie viele Container (in Millionen) in den Jahren 2009 bis 2017 im Hamburger Hafen be- und entladen wurden. Berechne die Spannweite und den Durchschnitt (das arithmetische Mittel) dieser Containerzahlen. Berechne, um wie viel Prozent die Containerzahl von 2010 auf 2011 gestiegen ist.
\par\smallskip\noindent \begin{tikzpicture}\begin{axis}[ybar,width=13.7cm,height=4.6cm,ymin=0,ymax=11.64,symbolic x coords={2009,2010,2011,2012,2013,2014,2015,2016,2017},xtick=data,enlarge x limits=0.08,bar width=6mm,ylabel={Container in Millionen},ylabel style={font=\small},tick label style={font=\small},nodes near coords,every node near coord/.append style={font=\scriptsize,/pgf/number format/.cd,use comma,1000 sep={}},ymajorgrids,grid style={mbgitter}]\addplot[fill=mbkasten,draw=black] coordinates {(2009,7.0) (2010,7.9) (2011,9.0) (2012,8.9) (2013,9.3) (2014,9.7) (2015,8.8) (2016,8.9) (2017,8.8)};\end{axis}\end{tikzpicture}\par
\rechenraster{4}
\fusshilfe{\mbox{19:~2,7 Mio.; 8,7 Mio.; $\tfrac{11}{79}$ $\approx$ 13,9 \%}}
\end{pfaufg}
\pfsteckt{steckt auch in Nr. 87 (P10 ’25)}
\pfstufekopf[smittelwert]{Mittelwert}{in 4 der letzten 5 Prüfungen}
\pfgruppe{}
\begin{pfaufg}{NI ’25}{20.}{}
Kira hat in 4 Handballspielen im Schnitt 8,5 Treffer erzielt. Wie viele Treffer sind das insgesamt?
\fusshilfe{\mbox{20:~4 · 8,5 = 34 Treffer}}
\end{pfaufg}
\begin{pfaufg}{P10 ’16}{21.}{1 BE}
Gib das arithmetische Mittel (den Durchschnitt) der Werte 8; 40; 60 an.
\rechenraster{1}
\fusshilfe{\mbox{21:~36}}
\end{pfaufg}
\pfgruppe{}
\begin{pfaufg}{}{22.}{}
In mehreren Spielen fielen diese Tore: $1$, $0$, $2$, $1$, $3$, $1$, $0$, $2$, $1$, $2$. Fülle die Häufigkeitstabelle aus. Berechne damit den Mittelwert.
\par\smallskip\noindent \sachtabelle{lc}{Tore & Anzahl}{0 & \leerzelle \\ 1 & \leerzelle \\ 2 & \leerzelle \\ 3 & \leerzelle}\par
\fusshilfe{\mbox{22:~(0 + 4 + 6 + 3) : 10; $1{,}3$}}
\end{pfaufg}
\begin{pfaufg}{BY ’22}{23.}{}
Die Spielerinnen einer Basketballmannschaft sind 154 cm, 170 cm, 162 cm, 158 cm und 156 cm groß. Berechne die durchschnittliche Körpergröße.
\par\smallskip\noindent \sachtabelle{l}{Name Tina, Danea, Pia, Lilly, Lotti}{Körpergröße 154, 170, 162, 158, 156 cm}\par
\fusshilfe{\mbox{23:~160 cm}}
\end{pfaufg}
\begin{pfaufg}{P10 ’20}{24.}{2 BE}
Ein Schwimmbad zählt in der ersten Woche der Sommerferien jeden Tag seine Besucher. Am Donnerstag kamen 180 Besucher. Im Balkendiagramm fehlt der Balken für Donnerstag noch. Berechne, wie viele Besucher in dieser Woche durchschnittlich an einem Tag im Schwimmbad waren.
\par\smallskip\noindent \begin{tikzpicture}\begin{axis}[xmin=0,xmax=400,xtick distance=100,tick label style={font=\small},xticklabel style={/pgf/number format/use comma,/pgf/number format/1000 sep={\,},/pgf/number format/fixed,/pgf/number format/precision=0},xmajorgrids,grid style={mbgitter},ytick={0,1,2,3,4,5,6},yticklabels={{Mo},{Di},{Mi},{Do},{Fr},{Sa},{So}},ymin=-0.5,ymax=6.5,minor x tick num=1,xminorgrids,nodes near coords,every node near coord/.append style={font=\scriptsize,/pgf/number format/use comma,/pgf/number format/1000 sep={\,},/pgf/number format/fixed,/pgf/number format/precision=0},bar width=4mm,xtick pos=bottom,ytick pos=left,xbar,y dir=reverse,width=11cm,height=6.8cm,xlabel={Anzahl der Besucher}]\addplot[fill=mbkasten,draw=black] coordinates {(170,0) (130,1) (210,2) (190,4) (240,5) (280,6)};\end{axis}\end{tikzpicture}\par
\pfzwfrage{Berechne die Summe der Besucher aller sieben Tage.}
\rechenraster{2}
\fusshilfe{\mbox{24:~200 Besucher}}
\end{pfaufg}
\begin{pfaufg}{}{25.}{}
Ein Busfahrer notiert die Fahrgäste bei acht Fahrten: $23$, $31$, $18$, $27$, $35$, $22$, $29$, $31$. Der Bus hat $32$ Plätze. Berechne Mittelwert und Spannweite. Hat der Bus immer genug Plätze?
\rechenraster{2}
\fusshilfe{\mbox{25:~35 - 18; $17$}}
\end{pfaufg}
\begin{pfaufg}{BW ’22}{26.}{}
Eine Messstation zwischen Russland und Alaska meldet diese Monatsmittel: Juli 10 °C, August 9 °C, September 5 °C, Oktober $-$1 °C, November $-$9 °C, Dezember $-$14 °C. Gib den Mittelwert an. Wie groß ist der Unterschied zwischen Juli und Dezember?
\par\smallskip\noindent \sachtabelle{l}{Monate Juli bis Dezember mit den Temperaturen}{}\par
\fusshilfe{\mbox{26:~Mittelwert 0 °C; Unterschied 24 K}}
\end{pfaufg}
\pfgruppe{}
\begin{pfaufg}{P10 ’17}{27.}{1 BE}
Gib den Durchschnitt (das arithmetische Mittel) dieser Längen an: 1,8 m; 1,7 m; 1,6 m; 1,7 m
\rechenraster{1}
\fusshilfe{\mbox{27:~1,7 m}}
\end{pfaufg}
\begin{pfaufg}{P10 ’21}{28.}{1 BE}
Jonas springt im Weitsprung 4,08 m, 3,88 m, 3,92 m und 4,12 m. Gib seine durchschnittliche Weite an.
\rechenraster{2}
\fusshilfe{\mbox{28:~4,00 m}}
\end{pfaufg}
\pfgruppe{}
\begin{pfaufg}{NI ’22}{29.}{}
Die Tabelle zeigt den mittleren Niederschlag in Hannover von Januar bis Juni. Berechne den Durchschnitt dieser sechs Monate.
\par\smallskip\noindent \sachtabelle{l}{Monat Januar bis Juni}{Niederschlag in mm: 52, 40, 48, 50, 63, 74}\par
\fusshilfe{\mbox{29:~327 mm : 6 = 54,5 mm}}
\end{pfaufg}
\begin{pfaufg}{}{30.}{}
Der Notenspiegel einer Klassenarbeit ist in der Tabelle. Berechne die Zahl der Schüler und den Notendurchschnitt.
\par\smallskip\noindent \sachtabelle{lcccccc}{Note & $1$ & $2$ & $3$ & $4$ & $5$ & $6$}{Anzahl & $4$ & $6$ & $7$ & $5$ & $2$ & $1$}\par
\par\noindent n = \leerfeld , x̄ = \leerfeld \par
\fusshilfe{\mbox{30:~$2{,}92$}}
\end{pfaufg}
\begin{pfaufg}{SH ’24}{31.}{}
Mittlere Sommertemperaturen in Deutschland: 2016: 17,8 °C; 2017: 17,9 °C; 2018: 19,3 °C; 2019: 19,2 °C; 2020: 18,2 °C. Der heißeste Sommer vor 1990 hatte etwa 18,5 °C. Ein Artikel sagt: So warm wie damals ist heute ein durchschnittlicher Sommer. Prüfe die Aussage.
\par\smallskip\noindent \sachtabelle{l}{Jahr 2016 bis 2020, mittlere Sommertemperatur in °C}{}\par
\fusshilfe{\mbox{31:~$\approx$ 18,5 °C – Aussage stimmt}}
\end{pfaufg}
\pfgruppe{}
\begin{pfaufg}{P10 ’15}{32.}{2 BE}
Eine Fußballmannschaft hatte in einer Saison 17 Heimspiele. Zu diesen Spielen kamen zusammen 680 353 Zuschauer. Berechne, wie viele Zuschauer im Durchschnitt bei einem Heimspiel waren. Runde sinnvoll.
\rechenraster{1}
\fusshilfe{\mbox{32:~680 $\tfrac{353}{17}$ $\approx$ 40 000 Zuschauer (40 021)}}
\end{pfaufg}
\pfgruppe{}
\begin{pfaufg}{NI ’22}{33.}{}
In einer 9. Klasse gaben 20 Schüler Noten: 3-mal die 1, 7-mal die 2, 5-mal die 3, 4-mal die 4, 1-mal die 5. Berechne die Durchschnittsnote. Die 10. Klasse hat 2,8. Welche Klasse hat besser bewertet?
\par\smallskip\noindent \sachtabelle{l}{Note 1–6}{Anzahl 3, 7, 5, 4, 1, 0\\ insgesamt 20}\par
\fusshilfe{\mbox{33:~53 : 20 = 2,65; Klasse 9 (2,65 < 2,8)}}
\end{pfaufg}
\pfgruppe{}
\begin{pfaufg}{P10 ’24}{34.}{3 BE}
Das Säulendiagramm zeigt, wie viel Niederschlag 2021 in Brandenburg in jedem Monat gefallen ist (in mm). Im ganzen Jahr 2021 fielen 581,7 mm Niederschlag. Zeige, dass im Durchschnitt pro Monat etwa 48,5 mm fielen. Ermittle, um wie viel Prozent der Niederschlag im April unter diesem Monatsdurchschnitt lag.
\par\smallskip\noindent \begin{tikzpicture}\begin{axis}[ymin=0,ymax=100,ytick distance=20,tick label style={font=\small},xtick pos=bottom,ytick pos=left,yticklabel style={/pgf/number format/use comma,/pgf/number format/1000 sep={\,},/pgf/number format/fixed,/pgf/number format/precision=1},ymajorgrids,grid style={mbgitter},xtick={0,1,2,3,4,5,6,7,8,9,10,11},xticklabels={{Jan},{Feb},{Mär},{Apr},{Mai},{Jun},{Jul},{Aug},{Sep},{Okt},{Nov},{Dez}},xmin=-0.5,xmax=11.5,bar width=5mm,nodes near coords,every node near coord/.append style={font=\scriptsize,/pgf/number format/use comma,/pgf/number format/1000 sep={\,},/pgf/number format/fixed,/pgf/number format/precision=1},ybar,width=13.5cm,height=5.2cm,ylabel={Niederschlag in mm},ylabel style={font=\small}]\addplot[fill=mbkasten,draw=black] coordinates {(0,52.2) (1,39.6) (2,35.3) (3,19.2) (4,49.3) (5,85.1) (6,36.1) (7,82.4) (8,21.7) (9,54.6) (10,64.1) (11,42.1)};\end{axis}\end{tikzpicture}\par
\pfzwfrage{Berechne den Unterschied zwischen dem Monatsdurchschnitt und dem Aprilwert.}
\rechenraster{3}
\fusshilfe{\mbox{34:~$\approx$ 48,5 mm; $\approx$ 60,4 \% unter dem Durchschnitt}}
\end{pfaufg}
\pfsteckt{steckt auch in Nr. 19 (P10 ’21); Nr. 14 (P10 ’22); Nr. 88 (P10 ’25)}
\pfstufekopf[swinkelberechnenundbe]{Winkel berechnen und beschriften}{in 3 der letzten 5 Prüfungen}
\pfgruppe{}
\begin{pfaufg}{}{35.}{}
Ein Satz lautet: „Ein Viertel der Befragten fährt Rad, das sind $25\,\%$.“ Kreuze an, was die Zahl $25$ hier ist.
\pfkreuzzeile{Anteil in \%}\pfkreuzzeile{Winkel in Grad}
\fusshilfe{\mbox{35:~Anteil in \%}}
\end{pfaufg}
\pfgruppe{}
\begin{pfaufg}{}{36.}{}
Der Streifen ist $10$ cm lang. Färbe $30\,\%$ des Streifens. Wie lang ist der gefärbte Teil?
\par\smallskip\noindent \streifenleer\par
\par\noindent \leerfeld[cm] \par
\fusshilfe{\mbox{36:~$3$ cm}}
\end{pfaufg}
\begin{pfaufg}{}{37.}{}
Ein Sektor im Kreisdiagramm steht für $10\,\%$. Wie groß ist sein Mittelpunktswinkel?
\par\noindent \leerfeld[°] \par
\fusshilfe{\mbox{37:~$36^\circ$}}
\end{pfaufg}
\begin{pfaufg}{NI ’26}{38.}{}
Männer nannten ihre Lieblingssportart: Fußball 25 \%, Basketball 15 \%. Berechne die Winkel im Kreisdiagramm.
\par\smallskip\noindent \sachtabelle{lrr}{Sportart & Anteil & Winkel}{Kreisdiagramm teils gezeichnet &  & }\par
\fusshilfe{\mbox{38:~Fußball 90°, Basketball 54°}}
\end{pfaufg}
\begin{pfaufg}{}{39.}{}
Zeichne in den Kreis einen Sektor mit $72^\circ$. Beginne am Radius. Beschrifte den Sektor mit „Bus“.
\par\smallskip\noindent \kreisleer\par
\rechenraster{1}
\fusshilfe{\mbox{39:~Sektor $72^\circ$ ab der Nullmarke, beschriftet}}
\end{pfaufg}
\pfgruppe{}
\begin{pfaufg}{BY ’22}{40.}{}
Eltern nutzen das Handy-Internet zu 60 \% für soziale Medien, zu 30 \% zum Surfen und zu 10 \% für Musik. Berechne die Mittelpunktswinkel für ein Kreisdiagramm.
\rechenraster{1}
\fusshilfe{\mbox{40:~216°, 108°, 36°}}
\end{pfaufg}
\begin{pfaufg}{P10 ’21}{41.}{2 BE}
Die Tabelle gibt an, welche Güter im Hamburger Hafen be- und entladen werden (Anteile in \%): Schrott 17, Flüssigkeiten ?, Container 66, Getreide 6, Sonstiges 1, gesamt 100. Ergänze den fehlenden Anteil für Flüssigkeiten. Schreib an jeden Sektor des Kreisdiagramms das passende Gut.
\par\smallskip\noindent \begin{tikzpicture}[line width=0.6pt,baseline=(current bounding box.north)]\begin{scope}\clip (0,0) -- (68.4:2.0) arc (68.4:90:2.0) -- cycle;\foreach \i in {-14,...,14}{\foreach \j in {-14,...,14}{\fill (\i*0.16,\j*0.16) circle (0.5pt);}}\end{scope}\draw (0,0) -- (90:2.0);\begin{scope}\clip (0,0) -- (32.4:2.0) arc (32.4:68.4:2.0) -- cycle;\foreach \i in {-12,...,12}{\draw[line width=0.3pt] (\i*0.18-2.5,-2.5) -- ++(5,5);}\end{scope}\draw (0,0) -- (68.4:2.0);\fill[black!60] (0,0) -- (-28.8:2.0) arc (-28.8:32.4:2.0) -- cycle;\draw (0,0) -- (32.4:2.0);\draw (0,0) -- (-28.8:2.0);\fill[black!10] (0,0) -- (-270:2.0) arc (-270:-32.4:2.0) -- cycle;\draw (0,0) -- (-32.4:2.0);\draw (0,0) circle (2.0); \fill (0,0) circle (1.3pt);\draw[line width=0.4pt] (79.2:1.40) -- (79.2:2.40);\draw[line width=0.4pt] ($(79.2:2.40)+(0,-0.00)$) -- ++(1.8,0);\draw[line width=0.4pt] (50.4:1.40) -- (50.4:2.40);\draw[line width=0.4pt] ($(50.4:2.40)+(0,-0.00)$) -- ++(1.8,0);\draw[line width=0.4pt] (1.8:1.40) -- (1.8:2.40);\draw[line width=0.4pt] ($(1.8:2.40)+(0,-0.00)$) -- ++(1.8,0);\draw[line width=0.4pt] (-30.6:1.40) -- (-30.6:2.40);\draw[line width=0.4pt] ($(-30.6:2.40)+(0,-0.00)$) -- ++(1.8,0);\draw[line width=0.4pt] (-151.2:1.40) -- (-151.2:2.40);\draw[line width=0.4pt] ($(-151.2:2.40)+(0,-0.00)$) -- ++(-1.8,0);\end{tikzpicture}\par
\rechenraster{1}
\fusshilfe{\mbox{41:~10 \%}}
\end{pfaufg}
\begin{pfaufg}{}{42.}{}
So kommen Kinder zur Schule: Rad $50\,\%$, Bus $40\,\%$, zu Fuß $10\,\%$. Berechne zuerst die Winkel. Zeichne dann das Kreisdiagramm.
\par\smallskip\noindent \kreisleer\par
\rechenraster{1}
\fusshilfe{\mbox{42:~180; $36$}}
\end{pfaufg}
\begin{pfaufg}{}{43.}{}
Ein Grundstück ist $800$ m² groß. Davon sind $160$ m² Haus, $400$ m² Rasen, $120$ m² Beete und $120$ m² Wege. Berechne die Anteile in Prozent und die Winkel für ein Kreisdiagramm. Welcher Teil ist am größten?
\rechenraster{1}
\fusshilfe{\mbox{43:~15\,\%; der Rasen}}
\end{pfaufg}
\begin{pfaufg}{nach BW ’21}{44.}{}
2019 wurden in Deutschland 2 Mio. t Erdöl gefördert: davon 1 Mio. t aus Schleswig-Holstein, 0,7 Mio. t aus Niedersachsen und 0,15 Mio. t aus Rheinland-Pfalz, Rest übrige Länder. Berechne die Winkel für ein Kreisdiagramm.
\rechenraster{2}
\fusshilfe{\mbox{44:~180°, 126°, 27°, 27°}}
\end{pfaufg}
\begin{pfaufg}{nach BW ’23}{45.}{}
Auf einem Glücksrad hat Grün die Wahrscheinlichkeit 30 \%, Gelb 1/5, Blau einen Sektor von 45°, der Rest ist rot. Berechne die Winkel aller Sektoren.
\rechenraster{2}
\fusshilfe{\mbox{45:~grün 108°, gelb 72°, blau 45°, rot 135°}}
\end{pfaufg}
\pfgruppe{}
\begin{pfaufg}{}{46.}{}
Im Streifen sind $45\,\%$ Rad und $30\,\%$ Bus eingetragen. Der Rest ist „zu Fuß“. Wie viel Prozent sind „zu Fuß“?
\par\noindent \leerfeld[\%] \par
\fusshilfe{\mbox{46:~$25\,\%$}}
\end{pfaufg}
\pfgruppe{}
\begin{pfaufg}{}{47.}{3 BE}
So kommen Kinder zur Schule: Rad $45\,\%$, Bus $30\,\%$, Auto $15\,\%$, zu Fuß $10\,\%$. Ordne die Sektoren im Kreisdiagramm den Anteilen zu und beschrifte sie.
\par\smallskip\noindent \kreisdiagrammleer{/45,/15,/30,/10}\par
\rechenraster{1}
\end{pfaufg}
\begin{pfaufg}{HH ’26}{48.}{}
Ein Glücksrad hat die Felder A bis E mit den Wahrscheinlichkeiten 1/16, 1/2, 1/8, 1/4, $\tfrac{1}{16}$. Berechne die Mittelpunktswinkel und zeichne das Glücksrad.
\par\smallskip\noindent \kreisleer[2]\par
\rechenraster{2}
\fusshilfe{\mbox{48:~A 22,5°; B 180°; C 45°; D 90°; E 22,5°}}
\end{pfaufg}
\begin{pfaufg}{P10 ’22}{49.}{3 BE}
750 Jugendliche wurden gefragt, wie oft sie Comics lesen. Ergebnis: mehrmals pro Woche 39 \%, einmal pro Woche 18 \%, einmal im Monat 27 \%, nie 16 \%. Im Kreisdiagramm ist nur der Sektor „nie“ beschriftet. Beschrifte die drei anderen Sektoren. Berechne den Mittelpunktswinkel für „nie“.
\par\smallskip\noindent \begin{tikzpicture}[line width=0.6pt,baseline=(current bounding box.north)]\draw (0,0) -- (90:2.0);\node[font=\small,anchor=241.2] at (61.2:2.15) {nie};\draw (0,0) -- (32.4:2.0);\draw (0,0) -- (-32.4:2.0);\draw (0,0) -- (-129.6:2.0);\draw (0,0) circle (2.0); \fill (0,0) circle (1.3pt);\draw[line width=0.4pt] (0:1.40) -- (0:2.40);\draw[line width=0.4pt] ($(0:2.40)+(0,-0.00)$) -- ++(1.8,0);\draw[line width=0.4pt] (-81:1.40) -- (-81:2.40);\draw[line width=0.4pt] ($(-81:2.40)+(0,-0.00)$) -- ++(1.8,0);\draw[line width=0.4pt] (-199.8:1.40) -- (-199.8:2.40);\draw[line width=0.4pt] ($(-199.8:2.40)+(0,-0.00)$) -- ++(-1.8,0);\end{tikzpicture}\par
\pfzwfrage{Berechne 16 \% von 360°.}
\rechenraster{2}
\fusshilfe{\mbox{49:~57,6°}}
\end{pfaufg}
\pfgruppe{}
\begin{pfaufg}{P10 ’24}{\llap{\textasteriskcentered\,}50.}{3 BE}
Familie Becker verbraucht am Tag etwa 508 Liter Leitungswasser: 1 Trinken und Kochen 12 l, 2 Körperpflege 256 l, 3 Toilettenspülung 136 l, 4 Wäsche waschen 60 l, 5 Gartenpflege 16 l, 6 Sonstiges 28 l. Das Kreisdiagramm zeigt diese Anteile; beschriftet sind nur „1 Trinken und Kochen“ und „5 Gartenpflege“. Beschrifte den Sektor für „Wäsche waschen“. Berechne den Mittelpunktswinkel für den Anteil „Trinken und Kochen“ am gesamten Verbrauch.
\par\smallskip\noindent \begin{tikzpicture}[line width=0.6pt,baseline=(current bounding box.north)]\draw (0,0) -- (90:2.0);\draw (0,0) -- (-91.4:2.0);\draw (0,0) -- (-111.2:2.0);\fill[black!25] (0,0) -- (-165:2.0) arc (-165:-153.7:2.0) -- cycle;\draw (0,0) -- (-153.7:2.0);\node[font=\small,anchor=20.65] at (-159.35:2.15) {5 Gartenpflege};\draw (0,0) -- (-165:2.0);\fill[black!25] (0,0) -- (-269.9:2.0) arc (-269.9:-261.4:2.0) -- cycle;\draw (0,0) -- (-261.4:2.0);\node[font=\small,anchor=274.35] at (-265.65:2.15) {1 Trinken und Kochen};\draw (0,0) circle (2.0); \fill (0,0) circle (1.3pt);\draw[line width=0.4pt] (-0.7:1.40) -- (-0.7:2.40);\draw[line width=0.4pt] ($(-0.7:2.40)+(0,-0.00)$) -- ++(1.8,0);\draw[line width=0.4pt] (-101.3:1.40) -- (-101.3:2.40);\draw[line width=0.4pt] ($(-101.3:2.40)+(0,-0.00)$) -- ++(-1.8,0);\draw[line width=0.4pt] (-132.45:1.40) -- (-132.45:2.40);\draw[line width=0.4pt] ($(-132.45:2.40)+(0,-0.00)$) -- ++(-1.8,0);\draw[line width=0.4pt] (-213.2:1.40) -- (-213.2:2.40);\draw[line width=0.4pt] ($(-213.2:2.40)+(0,-0.00)$) -- ++(-1.8,0);\end{tikzpicture}\par
\pfzwfrage{Berechne den Winkel für 60 l.}
\rechenraster{2}
\fusshilfe{\mbox{50:~$\approx$ 43°) beschriften; $\approx$ 8,5°}}
\end{pfaufg}
\pfgruppe{}
\begin{pfaufg}{P10 ’17}{51.}{3 BE}
Die Tabelle zeigt, als wievieltes Kind ihrer Mutter die Kinder des Jahrgangs 2013 in Deutschland zur Welt kamen (Anteil in Prozent): erstes Kind 49,4 \%, zweites Kind 34,4 \%, drittes Kind 11,2 \%, viertes oder weiteres Kind 5,0 \%. Im Kreisdiagramm ist einer dieser Anteile grau gefärbt. Schreib an den grauen Sektor, zu welchem Anteil er gehört. Berechne, wie groß der Mittelpunktswinkel für „drittes Kind“ sein muss. Zeichne diesen Sektor in den Kreis ein und beschrifte ihn.
\par\smallskip\noindent \begin{tikzpicture}[line width=0.6pt,baseline=(current bounding box.north)]\fill[black!25] (0,0) -- (-180:2.0) arc (-180:-162:2.0) -- cycle;\draw (0,0) -- (-180:2.0);\draw (0,0) circle (2.0); \fill (0,0) circle (1.3pt);\draw[line width=0.4pt] (-171:1.40) -- (-171:2.40);\draw[line width=0.4pt] ($(-171:2.40)+(0,-0.00)$) -- ++(-1.8,0);\draw[line width=0.4pt] ($(-171:2.40)+(0,-0.55)$) -- ++(-1.8,0);\draw[line width=0.4pt] ($(-171:2.40)+(0,-1.10)$) -- ++(-1.8,0);\end{tikzpicture}\par
\rechenraster{2}
\fusshilfe{\mbox{51:~$\approx$ 40°; Sektor 40° „drittes Kind“}}
\end{pfaufg}
\pfsteckt{steckt auch in Nr. 94 (P10 ’25)}
\pfstufekopf[sergnzen]{ergänzen}{in 2 der letzten 5 Prüfungen}
\pfgruppe{}
\begin{pfaufg}{}{52.}{}
Die y-Achse (senkrechte Achse) eines Säulendiagramms ist mit $0$, $20$, $40$, $60$, $80$ beschriftet. Kreuze an, ob die Achse bei null beginnt. Wenn nicht, trage den Startwert ein.
\pfkreuzzeile{ja, bei null}\pfkreuzzeile{nein, Startwert}
\fusshilfe{\mbox{52:~ja, bei null}}
\end{pfaufg}
\begin{pfaufg}{}{53.}{}
Sieh dir die y-Achse des Säulendiagramms an. Wie viel ist eine Hilfslinie wert? Gib nur diesen Wert an.
\par\smallskip\noindent \saeulenab[ymax=2,ystep=0.5,yfein=0.1,ylabel=Länge in m]{A/1.3,B/0.7,C/1.8}\par
\par\noindent \leerfeld[je]  Linie\par
\fusshilfe{\mbox{53:~$0{,}1$ je Linie}}
\end{pfaufg}
\pfgruppe{}
\begin{pfaufg}{VERA ’12}{54.}{}
Schüler zählen 500 Fahrzeuge: 300 Pkw, 120 Lkw, 60 Motorräder und Mofas, 20 Busse und andere. Zeichne die Anzahlen ins Stabdiagramm ein.
\par\smallskip\noindent \fbox{\parbox{0.9\linewidth}{\footnotesize Skizze: leeres Stabdiagramm mit vier Kategorien, Achse bis 300}}\par
\rechenraster{1}
\fusshilfe{\mbox{54:~Stäbe bei 300, 120, 60 und 20}}
\end{pfaufg}
\begin{pfaufg}{VERA ’11}{55.}{}
Julia hat die Tagesstrecken ihrer Radtour in ein Säulendiagramm gezeichnet. Den letzten Tag hat sie weggelassen; da fuhren sie nur 20 km. Ergänze die Säule.
\par\smallskip\noindent \fbox{\parbox{0.9\linewidth}{\footnotesize Skizze: Säulen, Tag 1–6: 68, 63, 58, 55, 45, 41 km; Tag 7 leer; Achse 0..80 km}}\par
\fusshilfe{\mbox{55:~Säule für Tag 7 bis 20 km}}
\end{pfaufg}
\begin{pfaufg}{}{56.}{}
An der y-Achse des Säulendiagramms stehen keine Zahlen. Die Säule A ist $6$ Kästchen hoch und zeigt $240$. Wie viel ist ein Kästchen wert?
\par\smallskip\noindent \saeulenab[ymax=400,ystep=40,yfein=40,ylabel=Anzahl,ohnezahlen]{A/240,B/320,C/120}\par
\par\noindent \leerfeld[je]  Kästchen\par
\fusshilfe{\mbox{56:~$40$ je Kästchen}}
\end{pfaufg}
\begin{pfaufg}{}{57.}{}
So kommen Kinder zur Schule: Rad $18$, Bus $24$, Auto $6$, zu Fuß $12$. Wähle eine passende Einteilung der y-Achse. Der größte Wert muss hineinpassen. Zeichne dann das Säulendiagramm.
\par\smallskip\noindent \saeulen{Rad/,Bus/,Auto/,zu Fuß/}{30}{6}{Kinder}\par
\rechenraster{1}
\fusshilfe{\mbox{57:~Säulen $18$, $24$, $6$, $12$}}
\end{pfaufg}
\begin{pfaufg}{}{58.}{}
Die Temperaturen der Monate in °C sind: Jan $1$, Feb $3$, Mär $7$, Apr $12$, Mai $16$. Zeichne den Verlauf als Liniendiagramm.
\par\smallskip\noindent \liniendia[ymax=20,ystep=5,yfein=1,ylabel=Temperatur in $^\circ$C]{Jan/,Feb/,Mär/,Apr/,Mai/}\par
\rechenraster{1}
\end{pfaufg}
\begin{pfaufg}{P10 ’20}{59.}{1 BE}
Ein Schwimmbad zählt in der ersten Woche der Sommerferien jeden Tag seine Besucher. Am Donnerstag kamen 180 Besucher. Im Balkendiagramm fehlt der Balken für Donnerstag noch. Zeichne im Diagramm den Balken für Donnerstag ein.
\par\smallskip\noindent \begin{tikzpicture}\begin{axis}[xmin=0,xmax=400,xtick distance=100,tick label style={font=\small},xticklabel style={/pgf/number format/use comma,/pgf/number format/1000 sep={\,},/pgf/number format/fixed,/pgf/number format/precision=0},xmajorgrids,grid style={mbgitter},ytick={0,1,2,3,4,5,6},yticklabels={{Mo},{Di},{Mi},{Do},{Fr},{Sa},{So}},ymin=-0.5,ymax=6.5,minor x tick num=1,xminorgrids,nodes near coords,every node near coord/.append style={font=\scriptsize,/pgf/number format/use comma,/pgf/number format/1000 sep={\,},/pgf/number format/fixed,/pgf/number format/precision=0},bar width=4mm,xtick pos=bottom,ytick pos=left,xbar,y dir=reverse,width=11cm,height=6.8cm,xlabel={Anzahl der Besucher}]\addplot[fill=mbkasten,draw=black] coordinates {(170,0) (130,1) (210,2) (190,4) (240,5) (280,6)};\end{axis}\end{tikzpicture}\par
\rechenraster{1}
\fusshilfe{\mbox{59:~Balken bei Do bis 180}}
\end{pfaufg}
\begin{pfaufg}{}{60.}{}
Ein Verein hat diese Mitgliederzahlen: Fußball $340$, Handball $280$, Tennis $120$, Turnen $200$. Welche Sportart zeigt die Säule mit dem Fragezeichen? Zeichne auch die Säule für Tennis ein.
\par\smallskip\noindent \saeulenab[ymax=400,ystep=100,yfein=20,ylabel=Mitglieder]{Fußball/340,?/280,Tennis/,Turnen/200}\par
\rechenraster{1}
\fusshilfe{\mbox{60:~280; $120$}}
\end{pfaufg}
\begin{pfaufg}{HH ’26}{61.}{}
Peter wirft eine Münze 25-mal: 12-mal Kopf, 13-mal Zahl. Gib die relativen Häufigkeiten als Bruch und in Prozent an.
\par\smallskip\noindent \sachtabelle{l}{Ergebnis Kopf/Zahl/Summe}{Spalten absolute Häufigkeit, relative Häufigkeit als Bruch, in Prozent}\par
\fusshilfe{\mbox{61:~Kopf $\tfrac{12}{25}$ = 48 \%, Zahl $\tfrac{13}{25}$ = 52 \%}}
\end{pfaufg}
\begin{pfaufg}{P10 ’18}{62.}{2 BE}
1000 Jugendliche haben ihr Lieblingsessen genannt: Schnitzel 80, Pizza 90, Kartoffelgerichte 160, Gemüsegerichte 160, Fischgerichte 110; die Zahlen für Pasta und Salat fehlen noch. Die fünf bekannten Zahlen sind in einem Balkendiagramm dargestellt. Die Achse „Anzahl“ hat aber noch keine Einteilung, und die Balken sind nur mit A bis E bezeichnet. Trag an der Achse eine passende Einteilung ein. Schreib zu jedem Balken das passende Gericht.
\par\smallskip\noindent \begin{tikzpicture}[x=0.45cm,y=0.45cm]\draw[mbgitter,line width=0.3pt,step=1] (0,0) grid (18,11);\draw[line width=0.7pt] (0,0) -- (0,11);\draw[line width=0.7pt,->] (0,0) -- (11,0);\fill[mbkasten,draw=black] (0,8.60) rectangle (5.5,9.40);\node[left,font=\small] at (0,9) {A};\draw[line width=0.4pt] (11,8.6) -- (17,8.6);\fill[mbkasten,draw=black] (0,6.60) rectangle (4.0,7.40);\node[left,font=\small] at (0,7) {B};\draw[line width=0.4pt] (11,6.6) -- (17,6.6);\fill[mbkasten,draw=black] (0,4.60) rectangle (4.5,5.40);\node[left,font=\small] at (0,5) {C};\draw[line width=0.4pt] (11,4.6) -- (17,4.6);\fill[mbkasten,draw=black] (0,2.60) rectangle (8.0,3.40);\node[left,font=\small] at (0,3) {D};\draw[line width=0.4pt] (11,2.6) -- (17,2.6);\fill[mbkasten,draw=black] (0,0.60) rectangle (8.0,1.40);\node[left,font=\small] at (0,1) {E};\draw[line width=0.4pt] (11,0.6) -- (17,0.6);\end{tikzpicture}\par
\pfzwfrage{Gib an, welche beiden Gerichte am häufigsten genannt wurden.}
\pfzwfrage{Berechne, wie viele Jugendliche ein Kästchen darstellt.}
\rechenraster{2}
\fusshilfe{\mbox{62:~20 je Gitterlinie}}
\end{pfaufg}
\pfgruppe{}
\begin{pfaufg}{}{63.}{}
Lies im Säulendiagramm ab, wie viele Besucher am Dienstag kamen.
\par\smallskip\noindent \saeulenab[ymax=80,ystep=10,yfein=10,ylabel=Besucher]{Mo/30,Di/40,Mi/70,Do/50}\par
\par\noindent \leerfeld[Besucher] \par
\fusshilfe{\mbox{63:~$40$ Besucher}}
\end{pfaufg}
\pfgruppe{}
\begin{pfaufg}{SH ’23}{64.}{}
Jonas würfelt 500-mal mit einem Würfel, dem eine Ecke fehlt. Augenzahl 1 bis 6, absolute Häufigkeit: 55; 58; 88; leer; leer; 110. Relative Häufigkeit: 0,110; leer; 0,176; 0,150; leer; 0,220. Ergänze die fehlenden Werte.
\par\smallskip\noindent \sachtabelle{l}{Spalten 1 bis 6 und Summe}{Zeile absolute Häufigkeit 55, 58, 88, leer, leer, 110, 500\\ Zeile relative Häufigkeit 0,110, leer, 0,176, 0,150, leer, 0,220, 1}\par
\fusshilfe{\mbox{64:~114; 0,228}}
\end{pfaufg}
\pfgruppe{}
\begin{pfaufg}{P10 ’23}{65.}{3 BE}
In einem Säulendiagramm sind zwei Werte aus der Tabelle schon als Säulen gezeichnet; die erste Säule gehört zu Spanisch. Trag an der y-Achse eine passende Einteilung ein. Schreib unter die zweite Säule, zu welcher Sprache sie gehört. Zeichne die Säule für Arabisch ein.
\par\smallskip\noindent \begin{tikzpicture}[x=0.45cm,y=0.45cm]\draw[mbgitter,line width=0.3pt] (0,0) grid (16,14);\draw[line width=0.7pt,->] (0,0) -- (0,14.6) node[above,font=\small] {Anzahl in Mio.};\draw[line width=0.7pt,->] (0,0) -- (16.6,0) node[right,font=\small] {Sprache};\fill[mbkasten,draw=black] (2,0) rectangle (4,7.8);\node[below=2pt,font=\small] at (3,0) {Spanisch};\fill[mbkasten,draw=black] (7,0) rectangle (9,10.5);\node[below=2pt,font=\small] at (8,0) {\leerzelle};\node[below=2pt,font=\small] at (13,0) {Arabisch};\end{tikzpicture}\par
\pfzwfrage{Zähle, wie viele Kästchen hoch die Säule für Spanisch ist.}
\fusshilfe{\mbox{65:~50 Mio.; Hindi; 5,8 Kästchen hoch}}
\end{pfaufg}
\pfstufekopf[saussageprfen]{Aussage prüfen}{in 1 der letzten 5 Prüfungen}
\pfgruppe{}
\begin{pfaufg}{}{66.}{}
Die y-Achse eines Säulendiagramms ist mit $0$, $50$, $100$, $150$ beschriftet. Kreuze an, wo sie beginnt.
\pfkreuzzeile{bei null}\pfkreuzzeile{nicht bei null}
\fusshilfe{\mbox{66:~bei null}}
\end{pfaufg}
\pfgruppe{}
\begin{pfaufg}{}{67.}{}
Laut einem Diagramm fahren $5\,\%$ der Schüler mit dem Roller. Mara sagt: „Jeder zwanzigste Schüler fährt Roller.“ Rechne nach und begründe, ob Mara recht hat.
\rechenraster{1}
\fusshilfe{\mbox{67:~die Aussage stimmt}}
\end{pfaufg}
\begin{pfaufg}{NI ’26}{68.}{}
Bei einer Wahl gab es 21 Stimmen. Amelie bekam 6, Sven 7. Haben die beiden zusammen mehr oder weniger als die Hälfte der Stimmen?
\rechenraster{1}
\fusshilfe{\mbox{68:~mehr als die Hälfte: 13 > 10,5}}
\end{pfaufg}
\pfgruppe{}
\begin{pfaufg}{}{69.}{}
Emma sagt: „Der Preis ist von Dienstag auf Mittwoch um mehr als $20$ Cent gestiegen.“ Begründe mit dem Säulendiagramm, ob Emma recht hat.
\par\smallskip\noindent \saeulenab[ymax=2,ystep=0.5,yfein=0.05,ylabel=Preis in €]{Mo/1.55,Di/1.6,Mi/1.75,Do/1.7}\par
\rechenraster{1}
\fusshilfe{\mbox{69:~Nein}}
\end{pfaufg}
\pfgruppe{}
\begin{pfaufg}{}{70.}{}
2022 kamen $240$ Besucher, 2023 kamen $320$ Besucher. Lina sagt: „2023 kamen um ein Drittel mehr Besucher als 2022.“ Rechne nach und begründe, ob Lina recht hat.
\rechenraster{1}
\fusshilfe{\mbox{70:~die Aussage stimmt}}
\end{pfaufg}
\begin{pfaufg}{}{71.}{}
Ein Verein verliert seit vier Jahren jedes Jahr $20$ Mitglieder. Die Zahl sank von $400$ auf $320$. Ina sagt: „In $16$ Jahren hat der Verein keine Mitglieder mehr.“ Begründe, ob Ina recht hat.
\rechenraster{1}
\end{pfaufg}
\begin{pfaufg}{nach P10 ’22}{72.}{2 BE}
Jedes Jahr werden $900$ Jugendliche befragt, ob sie täglich Sport treiben. Das Säulendiagramm zeigt, wie viele es sind. Jemand behauptet: „Von Jahr zu Jahr treiben immer weniger Jugendliche täglich Sport. Von 2020 bis 2023 hat sich die Anzahl um mehr als die Hälfte verringert.“ Prüfe beide Teile der Behauptung und begründe.
\par\smallskip\noindent \saeulenab[ymax=400,ystep=100,yfein=10,ylabel=Jugendliche]{2019/372,2020/348,2021/361,2022/330,2023/296}\par
\rechenraster{1}
\fusshilfe{\mbox{72:~$52$, weit unter der Hälfte}}
\end{pfaufg}
\pfgruppe{}
\begin{pfaufg}{P10 ’19}{\llap{\textasteriskcentered\,}73.}{2 BE}
2015 erschien ein Zeitungsartikel mit der Überschrift „Bücherlesen bald out?“. Darin heißt es: Die Buchhändler klagen, dass jedes Jahr dramatisch weniger Jugendliche Bücher lesen. Man fürchte, dass schon bald kein einziger Jugendlicher mehr zu einem Buch greift. Neben dem Text steht das Säulendiagramm. Schüler einer 10. Klasse halten zwei Aussagen des Artikels für falsch. Nenne eine davon und berichtige sie.
\par\smallskip\noindent \begin{tikzpicture}\begin{axis}[ymin=30,ymax=50,ytick distance=5,tick label style={font=\small},xtick pos=bottom,ytick pos=left,yticklabel style={/pgf/number format/use comma,/pgf/number format/1000 sep={\,},/pgf/number format/fixed,/pgf/number format/precision=0},ymajorgrids,grid style={mbgitter},xtick={0,1,2,3,4},xticklabels={{2011},{2012},{2013},{2014},{2015}},xmin=-0.5,xmax=4.5,bar width=7mm,nodes near coords,every node near coord/.append style={font=\scriptsize,/pgf/number format/use comma,/pgf/number format/1000 sep={\,},/pgf/number format/fixed,/pgf/number format/precision=0},ybar,width=8.0cm,height=5.2cm,ylabel={Bücher lesende Jugendliche in \%},ylabel style={font=\small}]\addplot[fill=mbkasten,draw=black] coordinates {(0,44) (1,42) (2,40) (3,39) (4,36)};\end{axis}\end{tikzpicture}\par
\rechenraster{1}
\fusshilfe{\mbox{73:~„extrem abnehmen“ falsch}}
\end{pfaufg}
\pfgruppe{}
\begin{pfaufg}{}{74.}{}
Eine Stadt zählt jedes Jahr die Fahrraddiebstähle: 2020 $410$, 2021 $380$, 2022 $395$, 2023 $350$. Die Zeitung schreibt: „Jedes Jahr gibt es weniger Diebstähle. 2030 gibt es keine mehr.“ Prüfe beide Teile der Aussage.
\rechenraster{1}
\end{pfaufg}
\begin{pfaufg}{NI ’23}{75.}{}
Notenspiegel: Note 1 2-mal, 2 4-mal, 3 6-mal, 4 5-mal, 5 2-mal, 6 1-mal. Lisa sagt: „Mehr als drei Viertel haben besser als 5 geschrieben.“ Hat sie recht? Begründe.
\par\smallskip\noindent \sachtabelle{l}{Note 1–6}{Anzahl 2, 4, 6, 5, 2, 1}\par
\fusshilfe{\mbox{75:~Ja. 17 von 20 = 85 \% > 75 \%.}}
\end{pfaufg}
\pfgruppe{}
\begin{pfaufg}{}{76.}{}
Die y-Achse beginnt bei $70$. Am Montag ist der Wert $72$, am Dienstag $74$. Ein Schüler sagt: „Die Säule für Dienstag ist doppelt so hoch, der Wert hat sich verdoppelt.“ Um wie viel Prozent ist der Wert wirklich gestiegen?
\rechenraster{1}
\fusshilfe{\mbox{76:~$\tfrac{1}{36} \approx 2{,}8\,\%$}}
\end{pfaufg}
\pfgruppe{}
\begin{pfaufg}{P10 ’22}{77.}{2 BE}
Von 2017 bis 2021 wurden jedes Jahr 1 200 Jugendliche gefragt, wie oft sie Bücher lesen. Das Säulendiagramm zeigt, wie viele von ihnen mehrmals pro Woche ein Buch lesen. Ben schaut auf das Diagramm und sagt: „Jedes Jahr lesen weniger Jugendliche Bücher. Von 2018 bis 2021 ist ihre Zahl um mehr als die Hälfte gesunken.“ Entscheide, ob Ben recht hat. Begründe deine Entscheidung.
\par\smallskip\noindent \begin{tikzpicture}\begin{axis}[ybar,width=8.5cm,height=4.6cm,ymin=390,ymax=508.75,symbolic x coords={2017,2018,2019,2020,2021},xtick=data,enlarge x limits=0.08,bar width=6mm,ylabel={Anzahl},ylabel style={font=\small},tick label style={font=\small},nodes near coords,every node near coord/.append style={font=\scriptsize,/pgf/number format/.cd,use comma,1000 sep={}},ymajorgrids,grid style={mbgitter}]\addplot[fill=mbkasten,draw=black] coordinates {(2017,468.0) (2018,432.0) (2019,485.0) (2020,470.0) (2021,400.0)};\end{axis}\end{tikzpicture}\par
\rechenraster{1}
\fusshilfe{\mbox{77:~2019 Anstieg 432 auf 485; $\approx$ 7,4 \%), nicht um die Hälfte}}
\end{pfaufg}
\pfgruppe{}
\begin{pfaufg}{P10 ’18}{\llap{\textasteriskcentered\,}78.}{4 BE}
Frauen wurden gefragt, wie oft sie Fleisch oder Wurst essen. Das Säulendiagramm zeigt die Anteile der Antworten. Prüfe die beiden Aussagen zum Säulendiagramm. Kreuze jeweils an, ob sie richtig, falsch oder nicht entscheidbar ist, und begründe. Aussage 1:
\par\smallskip\noindent \begingroup\renewcommand{\arraystretch}{1.5}\begin{tabular}{|p{5.5cm}|c|c|c|}\hline Aussage & richtig & falsch & nicht entscheidbar \\ \hline „Nicht einmal die Hälfte der Frauen isst 4- bis 6-mal pro Woche Fleisch oder Wurst.“ & $\square$ & $\square$ & $\square$ \\ \hline \multicolumn{4}{|l|}{\rule{0pt}{3.2ex}\footnotesize\color{mbgrau}Begründung:} \\ \hline „1- bis 3-mal pro Woche Fleisch oder Wurst isst jede 15. Frau.“ & $\square$ & $\square$ & $\square$ \\ \hline \multicolumn{4}{|l|}{\rule{0pt}{3.2ex}\footnotesize\color{mbgrau}Begründung:} \\ \hline\end{tabular}\endgroup\par
\pfzwfrage{Vergleiche 55 \% mit der Hälfte.}
\pfzwfrage{Schreibe „jede 15.“ als Prozentsatz.}
\fusshilfe{\mbox{78:~Aussage 1 falsch (55 \% > 50 \%); Aussage 2 falsch}}
\end{pfaufg}
\pfstufekopf[smedian]{Median}{in 1 der letzten 5 Prüfungen}
\pfgruppe{}
\begin{pfaufg}{P10 ’15}{79.}{1 BE}
Gib den Median (Zentralwert) dieser Messwerte an: 5 °C; 7 °C; 3 °C; 8 °C; 1 °C; 8 °C; 5 °C
\rechenraster{1}
\fusshilfe{\mbox{79:~5 °C}}
\end{pfaufg}
\pfgruppe{}
\begin{pfaufg}{P10 ’24}{80.}{2 BE}
Ermittle den Median (Zentralwert) dieser Datenreihe: 20 °C; 17 °C; 21 °C; 18 °C; 21 °C; 11 °C
\pfzwfrage{Ordne die Werte der Größe nach.}
\pfzwfrage{Berechne den Mittelwert der beiden mittleren Werte.}
\rechenraster{2}
\fusshilfe{\mbox{80:~19 °C}}
\end{pfaufg}

\par\Needspace*{0.55\textheight}\pfstufekopf[serkennen1]{Was ist gesucht?}{}
\begin{pfaufg}{}{81.}{}
Was ist gesucht? Kreuze an. Rechne nicht.
\par\smallskip\noindent \begingroup\renewcommand{\arraystretch}{1.4}\begin{tabular}{|>{\raggedright\arraybackslash}p{6.3cm}|>{\centering\arraybackslash\hyphenpenalty=0}p{1.55cm}|>{\centering\arraybackslash\hyphenpenalty=0}p{1.55cm}|>{\centering\arraybackslash\hyphenpenalty=0}p{1.55cm}|}\hline  & {\scriptsize\hspace{0pt}Mittelwert} & {\scriptsize\hspace{0pt}Minimum} & {\scriptsize\hspace{0pt}Median} \\ \hline Wie viele Besucher kamen im Durchschnitt pro Tag? & $\square$ & $\square$ & $\square$ \\ \hline Welche Sprache hat die meisten, welche die wenigsten Muttersprachler? & $\square$ & $\square$ & $\square$ \\ \hline Temperaturen: 20, 17, 21, 18, 21, 11 °C. Bestimme den Median. & $\square$ & $\square$ & $\square$ \\ \hline Liegt der Benzinpreis von Montag bis Freitag im Mittel unter 1,80 €? & $\square$ & $\square$ & $\square$ \\ \hline Wie groß ist der Unterschied zwischen dem nassesten und dem trockensten Monat 2021? & $\square$ & $\square$ & $\square$ \\ \hline Messwerte: 5, 7, 3, 8, 1, 8, 5 °C. Welcher Wert steht in der Mitte, wenn man sie ordnet? & $\square$ & $\square$ & $\square$ \\ \hline\end{tabular}\endgroup\par
\end{pfaufg}
\pfstufekopf[srckwrtsfehlenderwert]{rückwärts: fehlender Wert}{in 1 der letzten 5 Prüfungen}
\pfgruppe{}
\begin{pfaufg}{P10 ’22}{82.}{1 BE}
In einer Woche lag die Temperatur im Durchschnitt bei 20 °C. Ergänze die fehlende Temperatur am Samstag.
\par\smallskip\noindent \sachtabelle{lcccccc}{Mo & Di & Mi & Do & Fr & Sa & So}{21 °C & 20 °C & 19 °C & 22 °C & 20 °C & \leerzelle & 20 °C}\par
\pfzwfrage{Berechne, wie groß die Summe aller sieben Temperaturen sein muss.}
\fusshilfe{\mbox{82:~18 °C}}
\end{pfaufg}
\pfstufekopf[saussagenprfenauswirk]{Aussagen prüfen, Auswirkung erklären}{in 1 der letzten 5 Prüfungen}
\pfgruppe{}
\begin{pfaufg}{SH ’22}{83.}{}
Bei einem Punktespiel ergab sich: Durchschnitt 0,5 Punkte, Median 0 Punkte. Begründe, warum der Median hier wenig aussagt.
\rechenraster{1}
\end{pfaufg}
\pfgruppe{}
\begin{pfaufg}{nach P10 ’25}{84.}{3 BE}
Das Alter der Kinder einer Gruppe ist: $6$, $3$, $2$, $5$, $3$, $5$, $3$, $5$, $3$, $4$. Aussage 1 lautet: „Die Spannweite beträgt $4$ Jahre.“ Aussage 2 lautet: „Der Modalwert beträgt $5$ Jahre.“ Eine Aussage ist falsch. Kreuze sie an und berichtige sie.
\pfkreuzzeile{Aussage 1}\pfkreuzzeile{Aussage 2}
\fusshilfe{\mbox{84:~$4$}}
\end{pfaufg}
\pfgruppe{}
\begin{pfaufg}{}{85.}{}
Der Durchschnitt von fünf Werten ist $20$. Ein sechster Wert $32$ kommt dazu. Begründe ohne neue Rechnung, ob der Mittelwert steigt oder sinkt.
\rechenraster{1}
\end{pfaufg}
\begin{pfaufg}{}{86.}{}
Fünf Kinder bekommen so viel Taschengeld in €: $10$, $12$, $8$, $11$, $54$. Berechne Mittelwert und Median. Welcher Wert beschreibt das typische Taschengeld besser?
\rechenraster{1}
\fusshilfe{\mbox{86:~19; $11$}}
\end{pfaufg}
\begin{pfaufg}{P10 ’25}{87.}{2 BE}
In einer Kita-Gruppe sind die Kinder so alt (in Jahren): 5, 2, 1, 4, 2, 4, 2, 4, 2, 2. Eine der beiden Aussagen ist falsch: Kreuze die falsche Aussage an und berichtige sie.
\par\smallskip\noindent \begingroup\renewcommand{\arraystretch}{1.5}\begin{tabular}{|p{7.8cm}|c|p{3.4cm}|}\hline Aussage & falsch & Korrektur \\ \hline Die Spannweite beträgt 4 Jahre. & $\square$ &  \\ \hline Der Modalwert beträgt 4 Jahre. & $\square$ &  \\ \hline\end{tabular}\endgroup\par
\pfzwfrage{Zähle, wie oft jedes Alter vorkommt.}
\fusshilfe{\mbox{87:~4 stimmt}}
\end{pfaufg}
\pfgruppe{}
\begin{pfaufg}{P10 ’25}{\llap{\textasteriskcentered\,}88.}{3 BE}
In derselben Kita arbeiten 11 Erziehende. Sie sind so alt (in Jahren): 52, 57, 33, 44, 31, 30, 66, 58, 61, 41, 55. Bestimme das Durchschnittsalter der Erziehenden. Die älteste und die jüngste Person verlassen die Kita. Erkläre, wie sich dadurch das Durchschnittsalter verändert.
\pfzwfrage{Berechne die Summe aller 11 Alter.}
\rechenraster{3}
\fusshilfe{\mbox{88:~48 Jahre; 2 · 48}}
\end{pfaufg}
\pfstufekopf[sausprozentdarstellen]{aus Prozent darstellen}{in 1 der letzten 5 Prüfungen}
\pfgruppe{}
\begin{pfaufg}{}{89.}{}
Ein Sektor hat einen Mittelpunktswinkel von $72^\circ$. Wie viel Prozent des Kreises sind das?
\par\noindent \leerfeld[\%] \par
\fusshilfe{\mbox{89:~$20\,\%$}}
\end{pfaufg}
\pfgruppe{}
\begin{pfaufg}{NI ’26}{90.}{}
Familie Omari gibt 25 \% für Lebensmittel, 30 \% für Miete und 45 \% für Sonstiges aus. Stelle die Anteile in einem Streifendiagramm dar.
\par\smallskip\noindent \begin{tikzpicture}\draw[mbgitter,line width=0.3pt,step=0.5] (-0.5,-0.5) grid (10.5,2);\draw[line width=0.7pt] (0,0) rectangle (10,1.5);\end{tikzpicture}\par
\rechenraster{1}
\end{pfaufg}
\begin{pfaufg}{P10 ’18}{91.}{2 BE}
Frauen wurden gefragt, wie oft sie Fleisch oder Wurst essen. Das Säulendiagramm zeigt die Anteile der Antworten. Stelle die Anteile aus dem Säulendiagramm als Streifendiagramm dar. Der vorgezeichnete Streifen ist 10 cm lang.
\par\smallskip\noindent \begin{tikzpicture}\draw[mbgitter,line width=0.3pt,step=0.5] (-0.5,-0.5) grid (10.5,2);\draw[line width=0.7pt] (0,0) rectangle (10,1.5);\end{tikzpicture}\par
\pfzwfrage{Gib an, wie viele Zentimeter des Streifens 10 \% entsprechen.}
\rechenraster{2}
\end{pfaufg}
\pfgruppe{}
\begin{pfaufg}{}{92.}{}
Ein Kiosk verkauft Eis: im Juni Vanille $40\,\%$, Schoko $35\,\%$, Erdbeere $25\,\%$; im August Vanille $30\,\%$, Schoko $35\,\%$, Erdbeere $35\,\%$. Vergleiche die beiden Streifen: Welcher Anteil blieb gleich, welche änderten sich und um wie viele Prozentpunkte?
\rechenraster{1}
\end{pfaufg}
\begin{pfaufg}{NI ’24}{93.}{}
In Deutschland haben 41 \% Blutgruppe 0, 43 \% Blutgruppe A und 5 \% Blutgruppe AB. Der Rest hat Blutgruppe B. Ergänze den Anteil. Zeichne ein Streifendiagramm (10 cm = 100 \%).
\par\smallskip\noindent \sachtabelle{l}{Blutgruppe 0, A, B, AB}{Anteil 41 \%, 43 \%, leer, 5 \%}
\par\smallskip\noindent \begin{tikzpicture}\draw[mbgitter,line width=0.3pt,step=0.5] (-0.5,-0.5) grid (10.5,2);\draw[line width=0.7pt] (0,0) rectangle (10,1.5);\end{tikzpicture}\par
\fusshilfe{\mbox{93:~B: 11 \%}}
\end{pfaufg}
\pfgruppe{}
\begin{pfaufg}{P10 ’25}{94.}{3 BE}
In derselben Kita arbeiten 11 Erziehende. Sie sind so alt (in Jahren): 52, 57, 33, 44, 31, 30, 66, 58, 61, 41, 55. Ermittle, wie viel Prozent der Erziehenden jünger als 50 Jahre sind. Stelle diesen Anteil in einem Kreisdiagramm dar.
\par\smallskip\noindent \kreisleer[2.5]\par
\pfzwfrage{Zähle die Erziehenden unter 50 Jahren.}
\rechenraster{3}
\fusshilfe{\mbox{94:~$\approx$ 45,5 \%; $\approx$ 164°}}
\end{pfaufg}
\pfstufekopf[sfalscheneindruckerkl]{falschen Eindruck erklären}{in 1 der letzten 5 Prüfungen}
\pfgruppe{}
\begin{pfaufg}{nach P10 ’26}{\llap{\textasteriskcentered\,}95.}{}
Ein Diagramm zeigt die Temperatur an vier Tagen. Die y-Achse beginnt bei $18$ °C. Jemand sagt: „Am Mittwoch ($21$ °C) war es dreimal so warm wie am Montag ($19$ °C), das sieht man an den Säulen.“ Erkläre, wie dieser falsche Eindruck entsteht.
\rechenraster{1}
\fusshilfe{\mbox{95:~$2$ Grad, gut $10\,\%$}}
\end{pfaufg}
\begin{pfaufg}{}{96.}{}
Zwei Diagramme zeigen dieselben Werte $84$, $86$ und $88$. Im ersten beginnt die Achse bei $0$, im zweiten bei $82$. Erkläre, warum im zweiten Diagramm die dritte Säule dreimal so hoch wirkt wie die erste.
\rechenraster{1}
\end{pfaufg}
\pfgruppe{}
\begin{pfaufg}{}{97.}{}
Die y-Achse beginnt bei $90$. Die Werte sind $92$ und $96$. Tim sagt: „Die Säule ist doppelt so hoch, also hat sich der Wert verdoppelt.“ Finde den Fehler und schreibe die Aussage richtig.
\rechenraster{1}
\fusshilfe{\mbox{97:~$4$, gut $4\,\%$}}
\end{pfaufg}
\pfgruppe{}
\begin{pfaufg}{}{98.}{}
Ein Diagramm zeigt die Werte $1\,020$, $1\,040$ und $1\,060$. Seine Achse beginnt bei $1\,000$. Skizziere das Diagramm mit einer Achse ab $0$. Vergleiche, wie die beiden Diagramme wirken.
\par\smallskip\noindent \saeulen{2021/,2022/,2023/}{1200}{200}{Anzahl}\par
\rechenraster{1}
\end{pfaufg}
\pfgruppe{}
\begin{pfaufg}{P10 ’14}{\llap{\textasteriskcentered\,}99.}{2 BE}
Eine andere Bank zahlt nur 0,25 \% Zinsen im Jahr. Für ihre Werbung „Tolle Zinsen!“ zeichnet sie ein Säulendiagramm, in dem das Guthaben kräftig zu wachsen scheint. Erkläre, mit welchem Trick die Grafik die Entwicklung des Guthabens besser aussehen lässt, als sie ist.
\par\smallskip\noindent \begin{tikzpicture}\begin{axis}[ymin=1060,ymax=1072,ytick distance=2,tick label style={font=\small},xtick pos=bottom,ytick pos=left,yticklabel style={/pgf/number format/use comma,/pgf/number format/1000 sep={\,},/pgf/number format/fixed,/pgf/number format/precision=1},ymajorgrids,grid style={mbgitter},xtick={0,1,2},xticklabels={{nach einem Jahr},{nach zwei Jahren},{nach drei Jahren}},xmin=-0.5,xmax=2.5,bar width=7mm,nodes near coords,every node near coord/.append style={font=\scriptsize,/pgf/number format/use comma,/pgf/number format/1000 sep={\,},/pgf/number format/fixed,/pgf/number format/precision=1},ybar,width=8.2cm,height=5.2cm,ylabel={Guthaben in Euro},ylabel style={font=\small},xticklabel style={text width=2cm,align=center,font=\scriptsize}]\addplot[fill=mbkasten,draw=black] coordinates {(0,1064) (1,1067) (2,1069.5)};\end{axis}\end{tikzpicture}\par
\rechenraster{1}
\fusshilfe{\mbox{99:~y-Achse beginnt bei 1060 € statt bei 0}}
\end{pfaufg}
\pfsteckt{steckt auch in Nr. 73 (P10 ’19)}
\pfpruefstein{P10 ’26}
\begin{pfaufg}{}{}{}
Eine Tankstelle hat eine Woche lang jeden Tag um 18 Uhr den Benzinpreis notiert.
\par\smallskip\noindent \sachtabelle{lrrrrrrr}{Wochentag & Mo & Di & Mi & Do & Fr & Sa & So}{Benzinpreis in € pro Liter & 1,77 & 1,78 & 1,77 & 1,79 & 1,84 & 1,82 & 1,85}\par
\end{pfaufg}
\begin{pfaufg}{}{a)}{1 BE}
Gib die Spannweite der Benzinpreise in € pro Liter an.
\fusshilfe{\mbox{Prüfstein a):~0,08 € pro Liter}}
\end{pfaufg}
\begin{pfaufg}{}{b)}{2 BE}
Ein Kunde meint: Von Montag bis Freitag hat ein Liter Benzin an dieser Tankstelle im Durchschnitt weniger als 1,80 € gekostet. Zeige mit einer Rechnung, dass er recht hat.
\fusshilfe{\mbox{Prüfstein b):~1,79 € < 1,80 €}}
\end{pfaufg}
\begin{pfaufg}{}{c)}{2 BE}
An einer Tankstelle wurde eine Woche lang täglich um 18 Uhr der Benzinpreis notiert (€ pro Liter): Mo 1,77, Di 1,78, Mi 1,77, Do 1,79, Fr 1,84, Sa 1,82, So 1,85. Berechne die prozentuale Preissteigerung von Donnerstag auf Freitag.
\rechenplatz{3}
\fusshilfe{\mbox{Prüfstein c):~$\tfrac{5}{179}$ $\approx$ 2,8 \%}}
\end{pfaufg}
\begin{pfaufg}{}{d)}{1 BE}
Timo hat die Preise aus der Tabelle in ein Säulendiagramm übertragen. Die Säule für Sonntag fehlt noch. Zeichne sie ein.
\par\smallskip\noindent \begin{tikzpicture}\begin{axis}[ymin=1.75,ymax=1.85,ytick distance=0.02,tick label style={font=\small},xtick pos=bottom,ytick pos=left,yticklabel style={/pgf/number format/use comma,/pgf/number format/1000 sep={\,},/pgf/number format/fixed,/pgf/number format/precision=2},ymajorgrids,grid style={mbgitter},xtick={0,1,2,3,4,5,6},xticklabels={{Mo},{Di},{Mi},{Do},{Fr},{Sa},{So}},xmin=-0.5,xmax=6.5,bar width=7mm,minor y tick num=1,yminorgrids,ybar,width=10.2cm,height=5.2cm,ylabel={Benzinpreis in € pro Liter},ylabel style={font=\small}]\addplot[fill=mbkasten,draw=black] coordinates {(0,1.77) (1,1.78) (2,1.77) (3,1.79) (4,1.84) (5,1.82)};\end{axis}\end{tikzpicture}\par
\rechenplatz{3}
\fusshilfe{\mbox{Prüfstein d):~Säule So bis 1,85 (höchste Säule)}}
\end{pfaufg}
\begin{pfaufg}{}{*e)}{1 BE}
Timo schreibt in einer Bewertung der Tankstelle: „Mein Diagramm beweist: Von Donnerstag auf Freitag ist der Preis auf mehr als das Doppelte gestiegen.“ Begründe, wodurch Timos Diagramm diesen falschen Eindruck erweckt.
\par\smallskip\noindent \sachtabelle{lrrrrrrr}{Wochentag & Mo & Di & Mi & Do & Fr & Sa & So}{Benzinpreis in € pro Liter & 1,77 & 1,78 & 1,77 & 1,79 & 1,84 & 1,82 & 1,85}\par{\footnotesize\color{mbgrau}Säule So: 1,85}\par
\fusshilfe{\mbox{Prüfstein e):~Achse beginnt bei 1,75 € statt bei 0}}
\end{pfaufg}
\end{document}